% Page 047 — page imprimée 44
% Contenu seul : le préambule, l'en-tête et la pagination sont gérés par meyrueis.tex

Il passait régulièrement au presbytère, une ou deux fois par semaine, pour me
demander si j’étais libre. Il m’est arrivé assez souvent de l’accompagner. J’ai ainsi
pris contact, grâce à lui, avec une population catholique assez arriérée, de l’avis
même du docteur, mais qui s’est toujours montrée heureuse et même honorée d’être
visitée par un pasteur, dépourvu, de toute évidence, d’esprit de propagande ou de
prosélytisme.

A propos de ces traversées des Causses j’ai présenté à la mémoire un voyage en car
que j’y fis un jour, je ne sais plus pour quelle raison. Au départ de Meyrueis, au
dernier moment, un jeune curé vint s’asseoir à côté de moi, très exubérant et gentil
au possible, il engagea avec moi une conversation intéressante et animée. Au fur et à
mesure de l’entretien, une véritable sympathie s’établit entre nous deux. Quelle
agréable compagnie pour ce déplacement ! Puis, m’interrogeant plus directement, il
voulut savoir qui j’étais. Je n’avais aucune raison de le lui cacher : le pasteur de
Meyrueis. L’effet fut immédiat. Il me tourna carrément le dos et plus un mot ne sortit
de sa bouche. Quand il descendit , il ne me dit même pas au revoir !
Je n’ai jamais oublié cette attitude blessante, et il m’est arrivé, plus tard au cours de
réunions œcuméniques au Maroc, par exemple, avec d’autres curés, de raconter ce
petit incident pénible pour me réjouir qu’aujourd’hui, ce genre de comportement ne
serait ni possible ni concevable.

Même si des désaccords doctrinaux subsistent, au moins nous avons vis à vis les uns
des autres une attitude fraternelle grâce au Concile Vatican II et au pape Jean XXIII.

Enfin un mot sur les concerts. Il y en eut beaucoup au cours de notre séjour. Les plus
appréciés furent de deux sortes :

Ceux organisés avec le concours du célèbre violoncelliste de l’époque, d’origine
cévenole Jacques Serres. Il vient de mourir au moment où j’écris ces lignes. J’avais
béni, à Valleraugue son mariage avec la pianiste Ady Levastré. C’était vraiment du
grand art et le milieu estivant, très mélomane se trouvait à même de l’apprécier.
J’ajoute que le milieu paysan ne se montrait pas insensible au talent de cet artiste.

Je l’accompagnai à l’harmonium avec le souci évident de le mettre en valeur, mais en
m’efforçant quand même de dégager toutes les ressources de cet instrument, trop
décrié, de nos temples ou églises de montagne. Ensuite, de Meyrueis, nous étendîmes
notre audience à d’autres paroisses de la région, puis jusqu’à Pau, Millau, Montauban
et plus tard l’Afrique du nord

La seconde série de concerts se fit, en particulier, avec le concours d’une personne
qui faisait partie des vacanciers habituels, la Comtesse de Trémeuge. Cette femme
jeune, assez extraordinaire d’une grande beauté, avait reçu de la nature une voix
admirable. Elle n’appartenait pas à notre église, mais aimait fréquenter notre
presbytère. Elle admirait le dépouillement de la foi réformée. Devant cette artiste,
séduit par sa valeur et sa personne, je lui avais conseillé de lui permettre au moins du Bach du
Haendel ou du Franck. Je lui avais fait répéter, à la maison, des extraits de la Passion
selon St-Matthieu. Elle s’arrêta net et modestement, j’ajoute sagement, elle me
conseilla de lui permettre de chanter nos beaux cantiques, soigneusement
sélectionnés. Le succès fut total.
